

   var TARGET_SUM   = 90;   // approximation de 90 = 112 Mds du CAE - 21.9 (impulsion budgétaire primaire 2026) // https://www.ofce.sciences-po.fr/blog2024/fr/2026/20260228_MPPIMRS/ https://cae-eco.fr/static/pdf/Focus_124_Fipu_251016-final.pdf Je ne soustrais pas les mesures exceptionnelles (8.3G: mesures exceptionnelles 2026: IS GE, complémentaires santé) car j'ai ôté de la liste des mesures consensuelles pour un montant équivalent (mutualiser les achats de l'État: 3, lutte contre la fraude sociale et fiscale: 5, lutter contre prix de transferts: 4.6, prévenir chutes et rentes: 3.2, baisse des produits de santé 2.6, bonnes pratiques 2.8...)
	// Old: 112 Mds€ soit 3.7% de PIB, objectif 2025 (auquel il faudrait soustraire les 0.4% du budget 2025, on arriverait à 99Mds€)
    var ORANGE_MIN   = 76;    // Déficit primaire 2026 (3.2% - 0.72% d'impulsion non exceptionnelle 2026)*PIB (3042G
	// Old: seuil bas de la zone orange, correspond à un ajustement de 2.8% de PIB + les 0.4% déjà fait au budget 2025 = 3.2%, ce qui correspond au déficit primaire (donc si r = g ça stabilise la dette)
    var REPEAT_EVERY = 5;     // répétition des headers toutes les N lignes

   /* var labels = {
        1:  "Clarification des compétences afin de limiter les doublons entre échelons",
        2:  "Limiter l'usage des médicaments et les prescriptions abusives",
        3:  "Réduction des doublons de personnel associés à la clarification des compétences entre échelons",
        4:  "Mutualisation de la politique et des circuits d'achat de l'État",
        5:  "Supprimer l'exonération de TICPE sur le transport aérien et maritime",
        6:  "Hausse du barème de la taxe sur les billets d'avion",
        7:  "Taxes de 2 € sur les petits colis en provenance de pays tiers",
        8:  "Une mesure de taxation du capital (Zucman ou ISF)",
        9:  "Suppression de l'abattement fiscal des frais professionnels des retraités",
        10: "Augmenter l'impôt sur l'héritage (pacte Dutreil, assurances-vie, donation)",
        11: "Taux de TVA majoré de +5 % sur les produits de luxe",
        12: "Suppression du régime dérogatoire applicable au transport maritime",
        13: "Geler pendant 2 ans les dépenses de l'État au niveau de l'inflation (hors charge de la dette)",
        14: "Geler pendant 2 ans les dépenses des collectivités territoriales au niveau de l'inflation",
        15: "Diminution du crédit d'impôt recherche pour les grandes entreprises",
        16: "Réforme de l'hôpital : tarification, soins ambulatoires, nombre de lits",
        17: "Gel pendant 2 ans de l'indexation des retraites",
        18: "Diminuer les aides à l'apprentissage et exclure les diplômés au-delà de la licence",
        19: "Restaurer la taxe d'habitation sur les 20 % les plus aisés",
        20: "Supprimer les tarifs réduits sur le carburant",
        21: "Plafonner le quotient conjugal pour les couples aisés (à 10 000 €)",
        22: "Passer l'âge de l'ouverture des droits à 65 ans",
        23: "Suppression des allocations familiales et des APL au-dessus resp. de 5,2 SMIC & 39 000 €",
        24: "Supprimer les avantages fiscaux & sociaux applicables aux compléments de salaire",
        25: "Limiter la réduction des cotisations patronales",
        26: "Hausse de trois points du taux de prélèvement forfaitaire unique (PFU)",
        27: "Diminution d'un tiers des subventions à l'enseignement privé",
        28: "Augmentation de 4 mois de la durée de travail pour avoir droit au chômage",
        29: "Limiter les affections de longue durée",
        30: "Hausse de la CSG",
        31: "Restaurer la taxe d'habitation pour tous",
        32: "Gel des dépenses sociales au niveau de l'inflation",
        33: "Augmentation de la prise en charge par les complémentaires et les assurés",
        34: "Hausse du taux normal de TVA",
        35: "Suppression du taux réduit de TVA dans la restauration",
        36: "Hausse du taux d'imposition des sociétés à 33 %",
        37: "Diminution des exonérations patronales spécifiques à la taxe d'apprentissage",
        38: "Baisse des seuils de l'IR",
        39: "Suppression du taux réduit de TVA pour les travaux",
        40: "Suppression de MaPrimeRénov'",
        41: "Revenir sur l'augmentation des dépenses militaires à partir de 2024",
        42: "Hausse d'un point du taux de chaque tranche de l'impôt sur le revenu",
        43: "Suppression des exonérations des produits financiers (livret A, LDDS)",
        44: "Diminution globale des APL de 5 %",
        45: "Suppression du régime dérogatoire des taux en Outre-Mer",
        46: "Mécénat d'entreprises (dons des entreprises aux associations)",
        47: "Diminution des dépenses éducatives en fonction de la baisse démographique",
        48: "Suppression du quotient familial",
        49: "Suppression de l'accroissement des crédits en faveur de l'agriculture durable (dont plan protéines végétales et plan haies)",
        50: "Supprimer l'AME"
    };

    var amounts = {
        1:  6,      2:  2.6,   3:  3,     4:  1.5,   5:  4.1,
        6:  3.7,    7:  0.5,   8:  5,     9:  5,     10: 8,
        11: 3.5,    12: 1.4,   13: 13.8,  14: 7.6,   15: 1.5,
        16: 6,      17: 5.2,   18: 6.2,   19: 9.3,   20: 4.9,
        21: 1.1,    22: 10.6,  23: 1.3,   24: 18,    25: 11.4,
        26: 1.2,    27: 3.5,   28: 1.9,   29: 3.4,   30: 14.6,
        31: 21.8,   32: 1.3,   33: 1.9,   34: 6,     35: 4.2,
        36: 4.5,    37: 0.6,   38: 11.9,  39: 6.5,   40: 2.1,
        41: 12,     42: 6.8,   43: 2.2,   44: 0.8,   45: 2.2,
        46: 0.5,    47: 3.4,   48: 10.6,  49: 0.5,   50: 1.3
    };*/
   // --- ORGANISATION DE L'ÉTAT ---
	// Considered done: 4,2,7,46,49,6,45
	// Amounts done: 13.8=1.5+2.6+0.5+0.5+0.5+6+2.2 
	// Total (original): 259.5 / Total (new): 223.5
	// Excluded: 21 (21.8), 48 (10.6)
	// Modified: 8 (5 -> 11)
	
%%%%%%%%%%%%%%%%%%%%%%%%%%%%%%%%%%%%%%%%%%%%%%%%%%%%%%%%%%%%%%%%%%
Mesures et sources

	// Source de la plupart des mesures : CAE (https://claude.ai/public/artifacts/74cd0701-a4f0-499e-a409-cae1da524ee4)
	var policies = {
		// Baisse dépenses
		1:  { amount: 9,    label: "Éliminer les doublons entre échelons territoriaux"}, //1
		2: { amount: 21.4, label: "Geler pendant deux ans les dépenses de l'État et des collectivités territoriales"}, //13
		3: { amount: 12,   label: "Réduire les dépenses militaires (annuler la hausse prévue)"}, // 41
		//4: { amount: 2.1,  label: "Supprimer MaPrimeRénov' (subvention pour la rénovation des bâtiments)"}, // 40
		4: { amount: 1.5,  label: "Diminuer le crédit d'impôt recherche pour les grandes entreprises"}, // 15
		// Baisse prestations
		5: { amount: 4.3,  label: "Réduire les aides aux entreprises pour les contrats d'apprentissage"}, // 18
		// 6: { amount: 5.3,  label: "Réduire le remboursement de certains soins et des affections de longue durée les moins sévères"}, // 29
		6: { amount: 5.3,  label: "Réduire le remboursement de certains soins, des indemnités d'arrêt maladie, et du transport des patients"}, // 
		// Déployer à large échelle une politique de sécurisation des prescriptions et de responsabilisation de l'ensemble des acteurs 4.0
		//+Réformer les Affections de Longue Durée (ALD) : rendre imposables les indemnités journalières, introduire un ticket modérateur de 2%, recentrer les critères 3.4 (1.4 déjà réalisé PLF 2026, je compte 2)
		//+Augmenter la prise en charge par les complémentaires santé en relevant le ticket modérateur 1.5
		//+Augmenter la participation des assurés sur les médicaments à faible service médical, les cures thermales et les frais de transport hors affection de longue durée 0.4
		// Instaurer une participation des assurés aux dispositifs médicaux avec une franchise de 1 euro (plafond rehaussé à 50€) 0.3 (fait PLF 2026)
		//+Augmenter la participation des assurés aux dispositifs médicaux en relevant le ticket modérateur de 10 points 0.4
		//+Réduire la progression des dépenses de transport sanitaire en transférant certaines dépenses aux établissements de santé et en renforçant les contrôles 0.3
		// Assurer une régulation compatible avec la soutenabilité de notre système de santé 6
		//+Ajuster les dispositifs de prise en charge des indemnités journalières (franchise, durée, plafond) 2 (je compte 0.7 car 1.3 réalisé dans PLF 2026)
		7: { amount: 6.7,  label: "Retirer les aides sociales aux étrangers non-européens (RSA, allocations familiales, APL, minimum vieillesse)"}, // https://www.institutmontaigne.org/presidentielle-2022/eric-zemmour/retirer-les-aides-sociales-non-contributives-aux-etrangers-extra-europeens/ 
		8: { amount: 1.3,  label: "Supprimer l'Aide Médicale de l'État, qui couvre les soins urgents des étrangers en situation irrégulière"}, // 50
		9: { amount: 3.5,  label: "Diminuer d'un tiers des subventions à l'enseignement privé"}, // 27
		10: { amount: 3.4,  label: "Réduire les dépenses éducatives en proportion de la baisse démographique"}, // 47
		11: { amount: 1.9,  label: "Augmenter de 4 mois la durée de travail pour avoir droit au chômage"}, // 28
		12: { amount: 3.4,  label: "Geler les aides sociales (hors APL), diminuer les Aides Pour le Logement de 5% et les restreindre aux ménages modestes"}, // 32
		// Retraites
		13: { amount: 4.8,  label: "Réduire d'un demi-point la revalorisation des pensions de retraites pendant 4 ans"}, // 17 https://www.ipp.eu/publication/effets-budgetaires-et-redistributifs-des-mesures-socio-fiscales/ 8.9G, >1.4k: 0.52=1.9/3.6: https://www.senat.fr/rap/l25-131-1/l25-131-11.pdf seulement >2k: 0.41=0.7/1.7: https://www.ofce.fr/pdf/pbrief/2020/OFCEpbrief64.pdf
		14:  { amount: 4.4,    label: "Supprimer l'abattement de 10% sur les retraites pour le calcul de l'impôt sur le revenu"}, // 9 https://www.ipp.eu/wp-content/uploads/2025/06/chapitre_2_fiscalite_web.pdf
		15: { amount: 15, label: "Augmenter l'âge légal de départ à la retraite de 64 à 65 ans sauf pour les carrières longues ou avec pénibilité"}, // 22 https://www.ccomptes.fr/sites/default/files/2025-03/20250220-Situation-financiere-et-perspectives-du-systeme-de%20retraites.pdf https://www.institutmontaigne.org/presidentielle-2022/valerie-pecresse/repousser-lage-de-la-retraite-a-65-ans/
		// Hausse taxes
		17:  { amount: 14.2,  label: "Supprimer les exonérations de taxes (sur les carburants et autres) dans certains secteurs (maritime, aviation, routier, agriculture) et augmenter la taxe sur les billets d'avion"}, // 5
		// Supprimer le taux réduit de TVA pour les billets d'avion 0.3
		// Supprimer l'exonération de TICPE sur le transport maritime 0.6
		// Supprimer l'exonération de TICPE sur le transport aérien 3.5
		// Supprimer le tarif réduit pour le gazole routier de transport des marchandises 1.1
		// Supprimer le tarif réduit pour le gazole non routier non agricole 1.2
		// Supprimer le tarif réduit pour le gazole non routier agricole 1.4
		// Instaurer une éco-taxe poids lourds au kilomètre sur le réseau routier national hors autoroutes 1.0
		// Augmenter le barème de la taxe sur les billets d'avion 3.7
		// Supprimer la taxe au tonnage pour un retour à l'impôt sur les sociétés pour les compagnies maritimes 1.4
		18: { amount: 2.2,  label: "Soumettre les intérêts du Livret A et du LDDS à l'impôt sur le revenu et aux prélèvements sociaux"}, // 43
		16: { amount: 4.2,  label: "Aligner le taux de TVA dans la restauration (10%) sur le taux normal (20%)"}, // 35
		19: { amount: 6,    label: "Augmenter d'un point le taux normal de TVA (de 20% à 21%)"}, // 34
		20: { amount: 14.6, label: "Augmenter d'un point le taux de CSG (impôt qui s'applique à presque tous les revenus)"}, // 30
		21: { amount: 11.9, label: "Augmenter d'un point les taux de l'impôt sur le revenu et abaisser les seuils pour élargir le nombre de contribuables"}, // 38
		22: { amount: 11.4, label: "Augmenter les cotisations sur les salaires moyens (réduire les allègements de charges du CICE)"}, // 25
		23: { amount: 18,   label: "Supprimer les avantages fiscaux aux compléments de salaire (intéressement, participation, tickets-restaurants, etc.)"}, // 24
		24: { amount: 4.4,  label: "Augmenter le taux d'imposition des sociétés de 25% à 33,5%"}, // 36
		// Hausse taxes riches
		25: { amount: 9.3,  label: "Restaurer la taxe d'habitation sur les 20 % les plus aisés"}, // 19
		26:  { amount: 11,    label: "Rétablir une version renforcée de l'impôt sur la fortune des millionnaires (avec un taux de 2% pour les milliardaires)"}, // 8
		27: { amount: 9,    label: "Augmenter l'impôt sur l'héritage pour les 10% d'héritages les plus élevés"}, // 10
		28: { amount: 3.5,  label: "Créer un taux de TVA à 25 % pour les biens de luxe (montres, yachts, voitures de sport...)"}, // 11
		29: { amount: 1.2,  label: "Augmenter la taxe sur les revenus du capital de 30% à 33%"}, // 26
		30: { amount: 4.7,  label: "Augmenter l'impôt sur les revenus des plus aisés en rajoutant des tranches, avec un taux maximal à 65%"}, // https://www.institutmontaigne.org/presidentielle-2022/jean-luc-melenchon/rendre-limpot-sur-le-revenu-plus-progressif-avec-un-bareme-a-14-tranches/

%%%%%%%%%%%%%%%%%%%%%%%%%%%%%%%%%%%%%%%%%%%%%%%%%%%%%%%%%%%%%%%%%%

	/*
	x Hausse de deux points des cotisations retraites (mieux vaut augmenter CSG): 15.2, Cour https://www.ccomptes.fr/sites/default/files/2025-03/20250220-Situation-financiere-et-perspectives-du-systeme-de%20retraites.pdf
Geler prestations sociales: 7.9
x Supprimer le taux réduit de TVA pour travaux, hors rénovation énergétique: 4.5
TFF: 2G€ https://media.globalcitizen.org/14/40/14409e4b-b421-4785-880a-a6b27a02389b/rapport-sous-embargo-la-taxation-des-transactions-financieres-optimiser-le-dispositif-francais.pdf
		1:  { amount: 9,    label: "Éliminer les doublons entre échelons territoriaux"}, // 1
			  // Redéfinir qui fait quoi entre État, régions, départements et communes pour supprimer les chevauchements de compétences. Source chiffrage : Cour des comptes, rapport sur l'organisation territoriale (2023).
			
		//3:  { amount: 3,    label: "Réduction des doublons de personnel associés à la clarification des compétences entre échelons"},
			  // Corollaire RH de la mesure 1 : réduction des effectifs redondants une fois les compétences clarifiées. Source chiffrage : Cour des comptes (même rapport).
		//4:  { amount: 1.5,  label: "Mutualisation de la politique et des circuits d'achat de l'État"},
			  // Centraliser les achats publics (fournitures, prestations) pour obtenir de meilleurs prix via des marchés groupés. Source chiffrage : Cour des comptes, rapport sur les achats de l'État (2022).
		13: { amount: 21.4, label: "Geler pendant deux ans les dépenses de l'État et des collectivités territoriales"}, // 13.8+7.6 (policies 13, 14 des Progressistes)
			  // Maintenir les crédits ministériels en valeur constante (zéro croissance réelle) sur deux exercices. Source chiffrage : HCFP / trajectoire PLF 2025.
		//14: { amount: 7.6,  label: "Geler pendant deux ans les dépenses des collectivités territoriales"},
			  // Même principe appliqué aux budgets locaux (régions, départements, communes). Source chiffrage : HCFP / PLF 2025.
		41: { amount: 12,   label: "Réduire les dépenses militaires (annuler la hausse prévue)"},
			  // Annuler la trajectoire de remontée en puissance prévue par la LPM 2024-2030 (+3 Md€/an). Source chiffrage : LPM 2024-2030 (loi n°2023-703).
		// --- SANTÉ ---
		//2:  { amount: 2.6,  label: "Limiter l'usage des médicaments et les prescriptions abusives"},
			  // Réduire la surprescription et encourager les génériques/biosimilaires via des référentiels de bon usage. Source chiffrage : HCAAM, rapport sur la pertinence des soins (2022).
		//16: { amount: 6,    label: "Réforme de l'hôpital : tarification, soins ambulatoires, nombre de lits"},
			  // Refondre la T2A, développer la chirurgie ambulatoire et optimiser la capacité hospitalière. Source chiffrage : Cour des comptes / IGAS, rapports hospitaliers (2023).
		29: { amount: 5.3,  label: "Réduire le remboursement de certains soins et des affections de longue durée les moins sévères"}, // 5.3 (29) + 1.9 (33
			  // Resserrer les critères d'admission en ALD (exonération totale du ticket modérateur) pour les pathologies les moins sévères. Source chiffrage : HCAAM (2022).
		//33: { amount: 1.9,  label: "Augmentation de la prise en charge par les complémentaires et les assurés"},
			  // Transférer une part des remboursements de l'Assurance maladie vers les mutuelles et/ou augmenter le ticket modérateur. Source chiffrage : HCAAM.
		// --- RETRAITES ---
		17: { amount: 5.2,  label: "Gel des retraites pendant 2 ans"},
			  // Ne pas revaloriser les pensions sur l'inflation pendant deux ans (vs revalorisation annuelle au 1er janvier). Source chiffrage : DSS / COR.
		22: { amount: 10.6, label: "Augmenter l'âge légal de départ à la retraite de 64 à 65 ans"},
			  // Porter l'âge légal de départ à la retraite de 64 à 65 ans, au-delà de la réforme de 2023. Source chiffrage : COR, rapport annuel 2023.
		// --- PROTECTION SOCIALE ---
		//23: { amount: 1.3,  label: "Suppression des allocations familiales et des APL au-dessus resp. de 5,2 SMIC & 39 000 €"},
			  // Soumettre les allocations familiales et les APL à un plafond de revenus strict. Source chiffrage : CNAF / DHUP.
		28: { amount: 1.9,  label: "Augmentation de 4 mois de la durée de travail pour avoir droit au chômage"},
			  // Allonger la période de cotisation minimale ouvrant droit à l'assurance chômage. Source chiffrage : UNEDIC.
		32: { amount: 2.1,  label: "Geler les dépenses sociales hors APL et diminuer les APL de 5%"}, // 1.3 (32) + 0.8 (44) + 1.3 (23)
			  // Contenir l'ensemble des prestations sociales (hors retraites et santé) à croissance nulle en volume. Source chiffrage : HCFP.
		//44: { amount: 0.8,  label: "Diminution globale des APL de 5 %"},
			  // Réduction uniforme de 5 % du montant des Aides Personnalisées au Logement. Source chiffrage : DHUP / CNAF (dans le prolongement de la baisse de 5 € de 2017).
		50: { amount: 1.3,  label: "Supprimer l'Aide Médicale de l'État, qui couvre les soins urgents des étrangers en situation irrégulière"},
			  // Supprimer l'Aide Médicale de l'État, qui couvre les soins urgents des étrangers en situation irrégulière. Source chiffrage : DREES / PLF 2025 (~1,2 Md€ de dotation).
		// --- ÉDUCATION ---
		18: { amount: 6.8,  label: "Réduire les aides aux entreprises pour les contrats d'apprentissage et les supprimer pour les licenes et masters"}, // + 37
			  // Réduire les aides aux entreprises pour les contrats d'apprentissage et les supprimer pour les niveaux bac+3 et au-delà. Source chiffrage : Cour des comptes, rapport sur l'apprentissage (2023).
		27: { amount: 3.5,  label: "Diminuer d'un tiers des subventions à l'enseignement privé"},
			  // Réduire de 33 % les subventions de l'État aux établissements privés sous contrat. Source chiffrage : MEN (budget enseignement privé ~10 Md€).
		47: { amount: 3.4,  label: "Réduire les dépenses éducatives en proportion de la baisse démographique"},
			  // Adapter mécaniquement le nombre de postes et les dotations à la baisse du nombre d'élèves. Source chiffrage : MEN / Cour des comptes.
		// --- FISCALITÉ VERTE & TRANSPORT ---
		5:  { amount: 14.1,  label: "Supprimer les exonérations de taxes (principalement sur les carburants) dans certains secteurs (maritime, aviation, routiers, agriculture) et augmenter la taxe sur les billets d'avion"}, // 4.1 (5) + 1.4 (12) + 3.7 (6) + 4.9 (20)
			  // Mettre fin à l'exemption de taxe sur les carburants (kérosène, fioul lourd) dont bénéficient l'aviation et le shipping. Source chiffrage : rapport IGF sur les niches fiscales vertes (2023).
		//6:  { amount: 3.7,  label: "Hausse du barème de la taxe sur les billets d'avion"},
			  // Augmenter la taxe de solidarité sur les billets d'avion (TSBA), différenciée selon la classe et la destination. Source chiffrage : rapport Quinet / PLF discussions 2024.
		//7:  { amount: 0.5,  label: "Taxes de 2 € sur les petits colis en provenance de pays tiers"},
			  // Instaurer une taxe unitaire sur les envois e-commerce de faible valeur (type Shein, Temu) hors UE. Source chiffrage : propositions parlementaires 2024.
		//12: { amount: 1.4,  label: "Suppression du régime dérogatoire applicable au transport maritime"},
			  // Supprimer les avantages fiscaux spécifiques au secteur maritime (taxe au tonnage, exonérations sociales). Source chiffrage : rapport IGF sur les niches fiscales.
		//20: { amount: 4.9,  label: "Supprimer les tarifs réduits sur le carburant"},
			  // Supprimer les taux réduits de TICPE applicables à certains usages professionnels (agriculture, BTP, transports). Source chiffrage : rapport IGF sur les niches fiscales vertes.
		//45: { amount: 2.2,  label: "Aligner la TVA en Outre-Mer sur les taux en métropole"},
			  // Aligner la TVA et les accises ultramarines sur le droit commun métropolitain. Source chiffrage : rapport IGF sur les dépenses fiscales.
		//49: { amount: 0.5,  label: "Suppression de l'accroissement des crédits en faveur de l'agriculture durable (dont plan protéines végétales et plan haies)"},
			  // Annuler les hausses de crédits prévues pour la transition agricole dans le PLF 2025. Source chiffrage : PLF 2025.
		// --- FISCALITÉ DES ENTREPRISES ---
		15: { amount: 1.5,  label: "Diminuer le crédit d'impôt recherche pour les grandes entreprises"},
			  // Plafonner ou réduire le CIR pour les groupes au-dessus d'un certain CA, tout en le maintenant pour les PME. Source chiffrage : Cour des comptes, rapport sur le CIR (2021).
		24: { amount: 18,   label: "Supprimer les avantages fiscaux aux compléments de salaire (intéressement, participation, tickets-restaurants, etc.)"},
			  // Mettre fin aux exonérations de cotisations sur intéressement, participation, tickets-restaurants, etc. Source chiffrage : rapport IGAS/IGF sur les niches sociales (2021).
		25: { amount: 11.4, label: "Augmenter les cotisations sur les bas salaires (réduire les allègements de charges)"},
			  // Réduire l'ampleur ou le plafond des allègements Fillon sur les bas salaires. Source chiffrage : Cour des comptes / France Stratégie.
		36: { amount: 4.5,  label: "Augmenter le taux d'imposition des sociétés de 25% à 33%"},
			  // Revenir au taux d'IS de 33,1/3 % (avant sa baisse à 25 % entre 2018 et 2022). Source chiffrage : DGFiP.
		//37: { amount: 0.6,  label: "Diminution des exonérations patronales spécifiques à la taxe d'apprentissage"},
			  // Réduire les exemptions de taxe d'apprentissage dont bénéficient certains secteurs. Source chiffrage : DGFiP.
		//46: { amount: 0.5,  label: "Mécénat d'entreprises (dons des entreprises aux associations)"},
			  // Réduire ou plafonner la réduction d'IS au titre du mécénat d'entreprise (actuellement 60 % du don). Source chiffrage : Cour des comptes, rapport sur le mécénat (2018).
		// --- FISCALITÉ DES MÉNAGES ---
		8:  { amount: 5,    label: "Instaurer un impôt de 2% sur la fortune des milliardaires"},
		8:  { amount: 5,    label: "Restaurer l'impôt sur la fortune des millionnaires"},
			  // Instaurer un impôt minimum mondial sur les milliardaires (proposition Zucman/G20) ou rétablir l'ISF. Source chiffrage : EU Tax Observatory / Gabriel Zucman (2024).
		9:  { amount: 5,    label: "Restreindre l'abattement de 10% sur les retraites pour le calcul de l'impôt sur le revenu aux petites retraites"},
			  // Supprimer l'abattement de 10 % sur les pensions de retraite pour le calcul de l'IR. Source chiffrage : DGFiP / Cour des comptes.
		10: { amount: 9,    label: "Augmenter l'impôt sur l'héritage pour les 10% d'héritages les plus élevés"},
			  // Resserrer les dispositifs d'optimisation successorale (abattements Dutreil, clause bénéficiaire assurance-vie). Source chiffrage : France Stratégie / CAE.
		19: { amount: 9.3,  label: "Restaurer la taxe d'habitation sur les 20 % les plus aisés"},
			  // Rétablir la TH uniquement pour les ménages du dernier quintile de revenus, supprimée pour tous en 2023. Source chiffrage : DGFiP.
		//21: { amount: 1.1,  label: "Plafonner le quotient conjugal pour les couples aisés (à 10 000 €)"},
			  // Limiter l'avantage fiscal lié à l'imposition commune des couples à 10 000 € (vs ~30 000 € actuellement pour les revenus élevés). Source chiffrage : IPP.
		26: { amount: 1.2,  label: "Augmenter la taxe sur les revenus du capital de 30% à 33%"},
			  // Porter la flat tax sur les revenus du capital de 30 % à 33 %. Source chiffrage : DGFiP.
		30: { amount: 14.6, label: "Augmenter d'un point le taux de CSG (qui s'applique à tous les revenus)"},
			  // Augmenter le taux de Contribution Sociale Généralisée (1 point ≈ ~10 Md€). Source chiffrage : DSS / DGFiP.
		//31: { amount: 21.8, label: "Restaurer la taxe d'habitation pour tous"},
			  // Rétablir intégralement la TH supprimée entre 2018 et 2023 pour l'ensemble des résidences principales. Source chiffrage : DGFiP.
		38: { amount: 13, label: "Augmenter d'un point les taux de l'impôt sur le revenu et abaisser les seuils pour élargir le nombre de contribuables"}, // 11.9 (38) + 1.1 (21)
			  // Abaisser les seuils d'entrée dans chaque tranche du barème progressif, élargissant ainsi la base taxable. Source chiffrage : DGFiP.
		//42: { amount: 6.8,  label: "Hausse d'un point du taux de chaque tranche de l'impôt sur le revenu"},
			  // Augmenter tous les taux marginaux d'imposition de 1 point (ex. : 11 % → 12 %, 41 % → 42 %). Source chiffrage : DGFiP.
		43: { amount: 2.2,  label: "Soumettre les intérêts du Livret A et du LDDS à l'impôt sur le revenu et aux prélèvements sociaux"},
			  // Soumettre les intérêts du Livret A et du LDDS à l'IR et aux prélèvements sociaux. Source chiffrage : DGFiP.
		//48: { amount: 10.6, label: "Suppression du quotient familial"},
			  // Supprimer le mécanisme de parts fiscales liées aux enfants à charge dans le calcul de l'IR. Source chiffrage : IPP / DGFiP.
		// --- TVA ---
		11: { amount: 3.5,  label: "Créer un taux de TVA à 25 % pour les biens de luxe (montres, yachts, voitures de sport...)"},
			  // Créer un taux de TVA à 25 % pour les biens de luxe (montres, yachts, vêtements haut de gamme…). Source chiffrage : DGFiP / propositions parlementaires.
		34: { amount: 6,    label: "Augmenter d'un point du taux normal de TVA"},
			  // Augmenter le taux de TVA à 21 % (+1 pt). Source chiffrage : DGFiP (1 pt de TVA normale ≈ 6–7 Md€).
		35: { amount: 4.2,  label: "Aligner le taux de TVA dans la restauration (10%) sur le taux normal (20%)"},
			  // Revenir au taux normal pour la restauration (taux réduit à 10 % depuis 2012, coût ~3–4 Md€). Source chiffrage : rapport IGF.39: { amount: 6.5,  label: "Suppression du taux réduit de TVA pour les travaux"},
			  // Appliquer le taux normal aux travaux de rénovation dans les logements de plus de 2 ans (taux à 10 % actuellement). Source chiffrage : rapport IGF sur les dépenses fiscales.
		// --- LOGEMENT & DIVERS ---
		40: { amount: 2.1,  label: "Supprimer MaPrimeRénov'"},
			  // Supprimer le dispositif de subvention à la rénovation énergétique des logements (~2 Md€ de budget annuel). Source chiffrage : PLF 2025.
		*/